\documentclass[10pt,a4paper]{amsart}
\usepackage[margin=19mm]{geometry}
\usepackage{amsmath,amssymb,mathtools}
\usepackage[scr=rsfs,cal=euler]{mathalfa}
\usepackage{xfrac}
\usepackage[unicode,hidelinks,bookmarks=false]{hyperref}
\newcommand{\ie}{i.e.\ }
\newcommand{\cf}{cf.\ }
\newcommand{\resp}{respectively\ }
\newcommand{\Subsection}{\subsection}
\providecommand{\nobreakdash}{\nobreak\hbox{-}}
\setlength{\emergencystretch}{2em}
\multlinegap0pt
\usepackage{esint}
\usepackage{url}
\setcounter{MaxMatrixCols}{10}
\newtheorem{theorem}{Theorem}[section]
\newtheorem{acknowledgement}[theorem]{Acknowledgement}
\newtheorem{algorithm}[theorem]{Algorithm}
\newtheorem{axiom}[theorem]{Axiom}
\newtheorem{case}[theorem]{Case}
\newtheorem{claim}[theorem]{Claim}
\newtheorem{conclusion}[theorem]{Conclusion}
\newtheorem{condition}[theorem]{Condition}
\newtheorem{conjecture}[theorem]{Conjecture}
\newtheorem{corollary}[theorem]{Corollary}
\newtheorem{criterion}[theorem]{Criterion}
\newtheorem{definition}[theorem]{Definition}
\newtheorem{example}[theorem]{Example}
\newtheorem{exercise}[theorem]{Exercise}
\newtheorem{lemma}[theorem]{Lemma}
\newtheorem{notation}[theorem]{Notation}
\newtheorem{question}[theorem]{Question}
\newtheorem{proposition}[theorem]{Proposition}
\newtheorem{remark}[theorem]{Remark}
\newtheorem{solution}[theorem]{Solution}
\newtheorem{summary}[theorem]{Summary}
\numberwithin{equation}{section}  
\begin{document}
\title[Bottom spectrum and parabolicity of 3-manifolds with scalar ]{Bottom spectrum and parabolicity of 3-manifolds with scalar curvature lower bound}
\author{Ovidiu Munteanu; Jiaping Wang}
\date{}
\maketitle
\noindent\emph{Source and licence.} Bottom spectrum and parabolicity of 3-manifolds with scalar curvature lower bound. Version arXiv:2607.06508v1 (CC BY 4.0). This material is distributed under the Creative Commons Attribution 4.0 International licence, http://creativecommons.org/licenses/by/4.0/, as recorded in the arXiv OAI licence metadata for this exact version and shown on the abstract page https://arxiv.org/abs/2607.06508v1. Full rights notice: licenses/arxiv-2607.06508-CC-BY-4.0.txt.\par\smallskip
\noindent\textit{Editorial scope.} Counted unit: the original \texttt{theorem} environment \texttt{-} at source lines 455--465 together with its complete proof at lines 467--564; the delivered document keeps source lines 160--565 so that every referenced label and every citation resolves inside it. Native editable source is the paper's own concatenated TeX; the AMS wrapper is editorial formatting only. No publisher class, upstream script or unrelated artwork is redistributed.
\begin{equation}
P\left( \Omega \right) :=\int_{\partial ^{\ast }\Omega }u-\int_{N}\left(
\chi _{\Omega }-\chi _{\Omega _{0}}\right) h\,u \label{P}
\end{equation}
for all Caccioppoli subsets $\Omega\subset N$ with the symmetric difference
$\Omega\Delta \Omega_0$ contained in the interior of $N,$ where 
$\partial ^{\ast }\Omega$ is the reduced boundary of $\Omega$
given by $\partial \Omega\setminus \partial N$ and $\Omega_0$ is a fixed smooth domain in $N$
with $\partial_{+} N\subset \partial \Omega_0$ and $\partial_{-} N\cap \overline{\Omega_0}=\emptyset.$

When $n\leq 7,$ there exists a smooth domain $\Omega$ minimizing the functional $P$ (see \cite{Z, CL} for a proof). 
Such $\Omega$ is called a warped $\mu$-bubble following \cite{CL}. The existence proof shows that the reduced boundary
$\partial ^{\ast }\Omega$ of a warped $\mu$-bubble must be of positive distance away from both $\partial_{+} N$ and $\partial_{-} N.$

By the first variation formula,
the mean curvature $H$ of $\partial ^{\ast }\Omega$ is given by
\begin{equation}
H=h-\frac{u_{\nu }}{u},  \label{a2}
\end{equation}
where $\nu$ is the outward unit normal vector to $\Omega.$
Moreover, by the second variation formula, one has
\begin{eqnarray*} 
0 &\leq &\int_{\partial ^{\ast }\Omega }\left( \left\vert \nabla _{\partial ^{\ast }\Omega 
}\psi \right\vert ^{2}u-\left\vert A\right\vert ^{2}\psi ^{2}u-
\frac{u_{\nu }^{2}}{u}\psi ^{2}-\mathrm{Ric}(\nu, \nu)\,\psi
^{2}u\right) \\
&&+\int_{\partial ^{\ast }\Omega }\left( \left( \Delta u-\Delta _{\partial
^{\ast }\Omega }u\right) \psi ^{2}-H\,u_{\nu}\psi ^{2}-h_{\nu }\psi
^{2}u\right)
\end{eqnarray*}
for all $\psi \in C^{\infty }\left( \partial ^{\ast }\Omega \right),$ where 
$A$ and $\Delta _{\partial ^{\ast }\Omega }$ denote the second fundamental form
and the Laplacian of $\partial ^{\ast}\Omega,$ respectively.

In the case $n=3,$ by using a rearrangement idea due to Schoen and Yau \cite{SY}, 
the second variation formula can be rewritten into (see \cite{CL})
\begin{eqnarray}\label{a3}
0 &\leq &\int_{\partial ^{\ast }\Omega }\left( \left\vert \nabla _{\partial ^{\ast }\Omega 
}\psi \right\vert ^{2}u-\frac{1}{2}\left\vert A\right\vert ^{2}\psi ^{2}u-%
\frac{1}{2}\frac{u_{\nu }^{2}}{u}\psi ^{2}+K_{\partial ^{\ast }\Omega }\psi
^{2}u\right) \\
&&+\int_{\partial ^{\ast }\Omega }\left( \left( \Delta u-\Delta _{\partial
^{\ast }\Omega }u\right) \psi ^{2}-\frac{1}{2}h^{2}\psi ^{2}u-h_{\nu }\psi
^{2}u-\frac{1}{2}S\psi ^{2}u\right) \notag
\end{eqnarray}
for all $\psi \in C^{\infty }\left( \partial ^{\ast }\Omega \right),$ where $S$ is the scalar curvature of $M$ and
$K_{\partial^{\ast }\Omega }$ the Gauss curvature of $\partial ^{\ast}\Omega,$ respectively.

Now consider a complete noncompact three-dimensional manifold $M^{3}$ with one end and finite first Betti number. 
Let $d$ be a smoothing of the distance function to a fixed point $p\in M$ such that $d$ remains proper and satisfies 
$\left\vert \nabla d\right\vert \leq 2.$ 
For sufficiently large $L_0$ with all the representatives of the first homology $H_1(M)$ lying inside the set $\{d<L_0\},$
according to \cite{LT}, the boundary $\partial E_L$ of the only unbounded component $E_L$ of $M\setminus \{d\leq L\}$  
must be connected for all $L>L_0.$ 
In the following sections, we will consider warped $\mu$-bubbles in a $3$-dimensional manifold $N=E_{L_1}\setminus E_{L_2},$
where $L_1<L_2$ are regular values of $d$ with $L_0<L_1<L_2.$ Obviously,  
the set $N$ is a compact smooth manifold with boundary
components $\partial _{+}N=\partial E_{L_1}$ and $\partial _{-}N=\partial E_{L_2}.$
The function $h=h\left( d\right) $ will be taken as a smooth function depending on $d$ with
\begin{equation}
\lim_{d\rightarrow L_1}h\left( d\right) =+\infty ,\ \ \ \ \
\lim_{d\rightarrow L_2}h\left( d\right) =-\infty.  \label{a1}
\end{equation}

\begin{proposition}\label{X}
Let $M$ be a complete noncompact three dimensional manifold $M^{3}$ with one end and finite first Betti number
and $N=E_{L_1}\setminus E_{L_2}$ as above. Then for a warped $\mu$-bubble $\Omega$ in $N,$ there exists  
a connected component $\Sigma$ of $\partial ^{\ast }\Omega $ that separates the point
$p\in M$ from the infinity of $M.$ 
\end{proposition}

\begin{proof}
As $M$ has only one end, the set $E_{L_1}\setminus \Omega$ has one unbounded component $E_{\Omega}.$ 
Arguing as in \cite{LT}, one concludes that the smooth boundary of $E_{\Omega}$ must be connected. 
As noted above, the reduced boundary $\partial ^{\ast }\Omega $ is of positive distance from $\partial E_{L_1}.$
So $\partial E_{\Omega}$ must be a component of $\partial ^{\ast }\Omega.$ Obviously, it separates the point
$p\in M$ from the infinity of $M.$ 
This shows that $\Sigma=\partial E_{\Omega}$
is the desired component.
\end{proof}

\section{Bottom spectrum estimate}

We are now ready to prove the following sharp comparison result for the bottom spectrum. 

\begin{theorem}
Let $\left( M^{3},g\right) $ be complete noncompact three dimensional with
scalar curvature $S\geq -6K$ for $K\geq 0.$ Assume that $M$ has finitely
many ends and finite first Betti number. Then the bottom spectrum must satisfy
$\lambda _{1}\left( M\right)\leq K.$
\end{theorem}

\begin{proof}
Our argument is inspired by the work \cite{APX, HKKZ}.
Suppose by contradiction that there exists $\varepsilon >0$ such that 
\begin{equation}
\lambda _{1}\left( M\right) \geq K+\varepsilon.  \label{a0}
\end{equation}
Then there exists smooth function $w>0$ satisfying
\begin{equation}
\Delta w=-\lambda _{1}\left( M\right) w\text{ \ on \ }M.  \label{w}
\end{equation}

Since our argument uses only the geometry of $M$ outside an arbitrarily large compact
set, we may assume without loss of generality that $M$ has one end with
finite first Betti number.
As mentioned in the previous section, we will construct warped
$\mu$-bubbles in annulus region $N=E_{L_1}\setminus E_{L_2}$ of the end $E,$
where $$L_1=R-L \;\;\;\text{and}\;\;\; L_2=R+L,$$ with $R$ being an arbitrarily large number and $L$ a fixed number
to be chosen. Set $\Omega_0=E_{L_1}\setminus E_{R},$ $$u=w^{\gamma },$$ and 
\begin{equation*}
h\left( d\right) =\frac{2\pi }{\delta L}\tan \left( \frac{\pi }{2L}(R-d)\right)
\end{equation*} 
in (\ref{P}), where 
\begin{equation*}
\gamma =\frac{6}{2+\delta }
\end{equation*}
and 
\begin{equation*}
\delta =\frac{\varepsilon }{K+1}.
\end{equation*}
Without loss of generality, we may assume that $\varepsilon<\frac{1}{2}$, so  that $\delta <\frac{1}{2}$ as well. 
Obviously, the function $h=h\left( d\right) $ satisfies
\begin{equation*}
\lim_{d\rightarrow L_1}h\left( d\right) =+\infty ,\ \ \ \ \
\lim_{d\rightarrow L_2}h\left( d\right) =-\infty .  
\end{equation*}
So there exists a smooth warped $\mu$-bubble $\Omega.$ Moreover, by Proposition \ref{X},
there is a connected component $\Sigma \subset \partial ^{\ast }\Omega $
that separates $p\in M$ from the infinity of $M.$ On $\Sigma,$ the second variation formula (\ref{a3}) becomes
\begin{eqnarray*}
0 &\leq &\int_{\Sigma}\left( \left\vert \nabla _{\Sigma
}\psi \right\vert ^{2}u-\frac{1}{2}\left\vert A\right\vert ^{2}\psi ^{2}u-%
\frac{1}{2}\frac{u_{\nu }^{2}}{u}\psi ^{2}+K_{\Sigma}\psi
^{2}u\right) \\
&&+\int_{\Sigma }\left( \left( \Delta u-\Delta _{\Sigma}u\right) \psi ^{2}-\frac{1}{2}h^{2}\psi ^{2}u-h_{\nu }\psi
^{2}u-\frac{1}{2}S\psi ^{2}u\right)
\end{eqnarray*}%
for all $\psi \in C^{\infty }\left( \Sigma \right).$ 

For $\psi =\frac{1}{\sqrt{u}}$ it follows that%
\begin{eqnarray}
0 &\leq &\int_{\Sigma }\left( \frac{1}{4}\frac{\left\vert \nabla _{\Sigma
}u\right\vert ^{2}}{u^{2}}-\frac{1}{2}\frac{u_{\nu }^{2}}{u^{2}}+\frac{%
\Delta u-\Delta _{\Sigma }u}{u}\right)  \label{a3'} \\
&&+\int_{\Sigma }\left( -\frac{1}{2}\left\vert A\right\vert ^{2}+K_{\Sigma }-%
\frac{1}{2}h^{2}-h_{\nu }-\frac{1}{2}S\right) .  \notag
\end{eqnarray}%
Integrating by parts we have that 
\begin{equation*}
-\int_{\Sigma }\frac{\Delta _{\Sigma }u}{u}=\int_{\Sigma }\frac{\left\vert
\nabla _{\Sigma }u\right\vert ^{2}}{u^{2}}.
\end{equation*}%
Moreover, in view of (\ref{a2}), we have
\begin{equation*}
\left\vert A\right\vert ^{2}\geq \frac{1}{2}H^{2}=\frac{1}{2}\left( h-\frac{%
u_{\nu }}{u}\right) ^{2}.
\end{equation*}%
Therefore, (\ref{a3'}) becomes%
\begin{eqnarray}
0 &\leq &\int_{\Sigma }\left( -\frac{3}{4}\frac{\left\vert \nabla _{\Sigma
}u\right\vert ^{2}}{u^{2}}-\frac{1}{2}\frac{u_{\nu }^{2}}{u^{2}}-\frac{1}{4}%
\left( h-\frac{u_{\nu }}{u}\right) ^{2}\right)  \label{a4} \\
&&+\int_{\Sigma }\left( \frac{\Delta u}{u}+K_{\Sigma }-\frac{1}{2}%
h^{2}-h_{\nu }-\frac{1}{2}S\right) .  \notag
\end{eqnarray}%
Using the fact that $S\geq -6K$ together with the Gauss-Bonnet formula $\int_{\Sigma }K_{\Sigma }\leq 4\pi,$ we rewrite 
(\ref{a4}) into 
\begin{eqnarray}
0 &\leq &\int_{\Sigma }\left( -\frac{3}{4}\frac{\left\vert \nabla _{\Sigma
}u\right\vert ^{2}}{u^{2}}-\frac{3}{4}\frac{u_{\nu }^{2}}{u^{2}}+\frac{1}{2}h%
\frac{u_{\nu }}{u}\right)  \label{a5} \\
&&+\int_{\Sigma }\left( \frac{\Delta u}{u}-\frac{3}{4}h^{2}+\left\vert
\nabla h\right\vert +3K\right) +4\pi .  \notag
\end{eqnarray}%
As $\delta>0$, we have the inequality 
\begin{equation*}
\frac{1}{2}h\frac{u_{\nu }}{u}\leq \frac{3}{4}\frac{1}{1+\delta }h^{2}+\frac{%
1}{12}\left( 1+\delta \right) \frac{u_{\nu }^{2}}{u^{2}}.
\end{equation*}%
Plugging into (\ref{a5}) implies that
\begin{eqnarray*}
0 &\leq &\int_{\Sigma }\left( -\frac{3}{4}\frac{\left\vert \nabla _{\Sigma
}u\right\vert ^{2}}{u^{2}}-\frac{8-\delta }{12}\frac{u_{\nu }^{2}}{u^{2}}%
\right) \\
&&+\int_{\Sigma }\left( \frac{\Delta u}{u}-\frac{3}{4}\frac{\delta }{%
1+\delta }h^{2}+\left\vert \nabla h\right\vert +3K\right) +4\pi .
\end{eqnarray*}%
In particular, since $\delta<\frac{1}{2}$, 
this proves%
\begin{equation}
0\leq \int_{\Sigma }\left( -\frac{8-\delta }{12}\frac{\left\vert \nabla
u\right\vert ^{2}}{u^{2}}+\frac{\Delta u}{u}-\frac{1}{2}\delta
h^{2}+\left\vert \nabla h\right\vert +3K\right) +4\pi .  \label{a6}
\end{equation}%
Noting that
\begin{equation*}
h\left( d\right) =\frac{2\pi }{\delta L}\tan \left( \frac{\pi }{2L}(R-d)\right)
\end{equation*}%
and that $\left\vert \nabla d\right\vert \leq 2,$ one gets
\begin{equation*}
-\frac{1}{2}\delta h^{2}+\left\vert \nabla h\right\vert \leq \frac{2\pi ^{2}%
}{\delta L^{2}}.
\end{equation*}%
So (\ref{a6}) becomes
\begin{equation}
0\leq \int_{\Sigma }\left( -\frac{8-\delta }{12}\frac{\left\vert \nabla
u\right\vert ^{2}}{u^{2}}+\frac{\Delta u}{u}+3K+\frac{2\pi ^{2}}{\delta L^{2}%
}\right) +4\pi.  \label{a7}
\end{equation}
Since $u=w^{\gamma }$, we conclude from (\ref{a7}) that
\begin{equation}
0\leq \int_{\Sigma }\left( -\lambda _{1}\left( M\right) \gamma -\left(
\gamma -\frac{4+\delta }{12}\gamma ^{2}\right) \frac{\left\vert \nabla
w\right\vert ^{2}}{w^{2}}+3K+\frac{2\pi ^{2}}{\delta L^{2}}\right) +4\pi.
\label{a9}
\end{equation}%
Recalling that 
\begin{equation*}
\gamma =\frac{6}{2+\delta } \;\;\;\text{and}\;\;\;
\delta =\frac{\varepsilon }{K+1},
\end{equation*}
we have 
\begin{equation*}
-\left( \gamma -\frac{4+\delta }{12}\gamma ^{2}\right) =-\frac{3\delta }{%
\left( 2+\delta \right) ^{2}}\leq -\frac{\delta }{2},
\end{equation*}
and
\begin{equation*}
-\lambda _{1}\left( M\right) \gamma +3K\leq -\varepsilon.
\end{equation*}
Hence, (\ref{a9}) becomes 
\begin{equation*}
0\leq \int_{\Sigma }\left( -\frac{\varepsilon }{2\left( K+1\right) }\frac{%
\left\vert \nabla w\right\vert ^{2}}{w^{2}}-\varepsilon +\frac{2\left(
K+1\right) \pi ^{2}}{\varepsilon L^{2}}\right) +4\pi .
\end{equation*}%
By setting
\begin{equation*}
L=\frac{2\left( K+1\right) \pi }{\varepsilon }
\end{equation*}%
we obtain that 
\begin{equation*}
\int_{\Sigma }\left( \frac{\left\vert \nabla w\right\vert ^{2}}{w^{2}}%
+1\right) \leq \frac{8\left( K+1\right) \pi }{\varepsilon }.
\end{equation*}

In conclusion, for each large $R>1,$ there exists a compact surface $\Sigma_{R},$ 

\begin{equation*}
\Sigma _{R}\subset B_{p}\left( R+2L\right) \setminus B_{p}\left( R-2L\right),
\end{equation*}
that separates the fixed point $p\in M$ from the infinity of $M$ and satisfies

\begin{equation}
\int_{\Sigma _{R}}\left( \frac{\left\vert \nabla w\right\vert ^{2}}{w^{2}}%
+1\right) \leq \frac{8\left( K+1\right) \pi }{\varepsilon }.  \label{a10}
\end{equation}
To complete our argument, let $N_{R}$ be the unbounded component of $M\setminus \Sigma _{R}$ and 
$D_{R}:=M\setminus N_{R}$. That is, $D_{R}$ is the bounded domain of $M$ with
boundary $\Sigma _{R}.$ Clearly,  as $L$ is a fixed constant, there exists a sequence of $R_i\to \infty$
such that $D_{R_i}$ is increasing in $i$ and
\begin{equation}
M=\cup _{i=1}^\infty D_{R_i}.  \label{a11}
\end{equation}

By (\ref{w}), we have
\begin{equation*}
\lambda _{1}\left( M\right) =-\Delta \ln w-\left\vert \nabla \ln
w\right\vert ^{2}\leq -\Delta \ln w.
\end{equation*}%
Integrating the equation over $D_R,$ we conclude that 
\begin{eqnarray*}
\lambda _{1}\left( M\right) \mathrm{Vol}\left( D_{R}\right) &\leq
&-\int_{D_{R}}\Delta \ln w \\
&\leq &\int_{\Sigma _{R}}\left\vert \nabla \ln w\right\vert \\
&\leq &\frac{1}{2}\int_{\Sigma _{R}}\left( \left\vert \nabla \ln
w\right\vert ^{2}+1\right) \\
&\leq &\frac{4\left( K+1\right) \pi }{\varepsilon },
\end{eqnarray*}
where we have used (\ref{a10}).
Hence, the volume of $D_{R}$ is uniformly bounded from above independent of $R.$ In view of (\ref{a11}), it shows that
$M$ must have finite volume as well. This is an obvious contradiction to the fact that the bottom spectrum of $M$ is positive, see Chapter 22 in \cite{Li}.
\end{proof}

In particular, we have the following consequence.

\begin{corollary}
Let $\left( M^{3},g\right) $ be complete noncompact three dimensional with
scalar curvature $S\geq 0$. Assume that $M$ has finitely many ends and
finite first Betti number. Then $\lambda _{1}\left( M\right) =0$.
\end{corollary}

The following result strengthens the above corollary.


\begin{theorem}
Let $\left( M^{3},g\right) $ be a complete noncompact three dimensional manifold with
scalar curvature $S\geq 0.$ Assume that $M$ has finitely many ends and
finite first Betti number. Then there exists a sequence of geodesic balls $B_p(R_i)$ with
$R_{i}\rightarrow \infty $ such that their first Dirichlet eigenvalues satisfy
 
\begin{equation*}
\lambda _{1}\left( B_{p}\left( R_{i}\right) \right) \leq \frac{C}{R_{i}^{2}}
\end{equation*}%
for a universal constant $C>0.$
\end{theorem}

\begin{proof}
We argue by contradiction. Suppose that
\begin{equation}
\lambda _{1}\left( B_{p}\left( R\right) \right) \geq \frac{A}{R^{2}}
\label{m1}
\end{equation}
for all $R>R_{0},$ where both $R_0$ and $A>1$ are constants as large as one wants to specify.

Again, without loss of generality, we may assume that $M$ has one end and its first
Betti number is zero. As in the previous theorem, 
we will construct warped
$\mu$-bubbles in annulus region $N=E_{L_1}\setminus E_{L_2}$ of the end $E,$
where $L_1=R-L$ and $L_2=R+L$ with $R$ being an arbitrarily large number and $L=\frac{R}{4}.$ 
In (\ref{P}), set $\Omega_0=E_{L_1}\setminus E_{R},$ $h$ the function given by
\begin{equation*}
h\left( d\right) =\frac{2\pi }{L}\tan \left( \frac{\pi }{2L}(R-d)\right)
\end{equation*} 
and $u$ the positive Dirichlet eigenfunction of $B_{p}\left(2R\right),$ that is,
\begin{equation*}
\Delta u=-\lambda _{1}\left( B_{p}\left( 2R\right) \right) u.
\end{equation*}
Then there exists a smooth warped $\mu$-bubble $\Omega$ in $N.$ Moreover, by Proposition \ref{X},
there is a connected component $\Sigma \subset \partial ^{\ast }\Omega $
that separates $p\in M$ from the infinity of $M.$

Now the second variation formula (\ref{a4}) implies that
\begin{equation*}
0\leq \int_{\Sigma }\left( \frac{\Delta u}{u}-\frac{1}{2}\frac{\left\vert
\nabla u\right\vert ^{2}}{u^{2}}+K_{\Sigma }-\frac{1}{2}h^{2}-h_{\nu }-\frac{%
1}{2}S\right) .
\end{equation*}%
Since $S\geq 0$ and $\int_{\Sigma }K_{\Sigma }\leq 4\pi $, we get 
\begin{equation}
\int_{\Sigma }\left( -\frac{\Delta u}{u}+\frac{1}{2}\frac{\left\vert \nabla
u\right\vert ^{2}}{u^{2}}+\frac{1}{2}h^{2}-\left\vert \nabla h\right\vert
\right) \leq 4\pi .  \label{m2}
\end{equation}%
Noting that $\Sigma \subset B_{p}\left( 2R\right)\setminus B_{p}\left( \frac{R}{2}\right) $ and that
$h$ satisfies
\begin{equation*}
\frac{1}{2}h^{2}-\left\vert \nabla h\right\vert \geq -\frac{C}{R^{2}},
\end{equation*}%
we conclude from (\ref{m2}) that 
\begin{equation*}
\int_{\Sigma }\left( \lambda _{1}\left( B_{p}\left( 2R\right) \right) -\frac{%
C}{R^{2}}+\frac{1}{2}\frac{\left\vert \nabla u\right\vert ^{2}}{u^{2}}%
\right) \leq 4\pi .
\end{equation*}%
By (\ref{m1}) it follows that 
\begin{equation}
\int_{\Sigma }\frac{\left\vert \nabla u\right\vert }{u}\leq CR.  \label{m3}
\end{equation}%
Denoting with $D$ the compact domain in $M$ with $\partial D=\Sigma,$ we have
\begin{eqnarray*}
\lambda _{1}\left( B_{p}\left( 2R\right) \right) \mathrm{Vol}\left( D\right) 
&=&\int_{D}\left( -\Delta \log u-\left\vert \nabla \log u\right\vert
^{2}\right)  \\
&\leq &-\int_{D}\Delta \log u \\
&\leq &\int_{\Sigma }\frac{\left\vert \nabla u\right\vert }{u} \\
&\leq &CR.
\end{eqnarray*}%
This shows that 
\begin{equation}
\lambda _{1}\left( B_{p}\left( 2R\right) \right) V_{p}\left( \frac{R}{2}%
\right) \leq CR  \label{m4}
\end{equation}%
for all $R>R_{0},$ where $V_p(r)$ denotes the volume of the geodesic ball $B_p(r).$
Appealing to (\ref{m1}), we conclude that 
\begin{equation}
V_{p}\left( R\right) \leq \frac{C_{0}}{A}R^{3}  \label{m5}
\end{equation}%
for all $R>R_{0},$ where $C_{0}>0$ is a universal constant. Let $\phi$ be 
the cut-off function satisfying $\phi =1$ on $B_{p}\left( R\right),$ $\phi =0$
on $M\setminus B_{p}\left( 2R\right) $ and $\left\vert \nabla \phi
\right\vert \leq \frac{2}{R}.$ Then
\begin{eqnarray*}
\lambda _{1}\left( B_{p}\left( 2R\right) \right) V_{p}\left( R\right)  &\leq
&\lambda _{1}\left( B_{p}\left( 2R\right) \right) \int_{M}\phi ^{2} \\
&\leq &\int_{M}\left\vert \nabla \phi \right\vert ^{2} \\
&\leq &\frac{4}{R^{2}}V_{p}\left( 2R\right).
\end{eqnarray*}%
Using (\ref{m1}) and (\ref{m5}) we conclude that 
\begin{equation*}
\frac{A}{(2R)^{2}}V_{p}\left( R\right) \leq \lambda _{1}\left( B_{p}\left(
2R\right) \right) V_{p}\left( R\right) \leq \frac{4}{R^{2}}V_{p}\left(
2R\right) \leq \frac{4}{R^{2}}\frac{C_{0}}{A}\left( 2R\right) ^{3}.
\end{equation*}%
Therefore,
\begin{equation*}
V_{p}\left( R\right) \leq \frac{2^7 C_{0}}{A^{2}}R^{3}
\end{equation*}%
for all $R>R_{0}.$ Iterating this argument $m$ times, we get
that 
\begin{equation*}
V_{p}\left( R\right) \leq \frac{2^{7m}C_{0}}{A^{m}}R^{3}
\end{equation*}%
for all $R>R_{0}.$ Letting $m\rightarrow \infty $ we arrive at an obvious contradiction if $A>2^7$.
\end{proof}


\begin{thebibliography}{99}
\bibitem[APX]{APX} G. Antonelli, M. Pozzetta, and K. Xu, A sharp splitting theorem, arXiv:2412.12707.
\bibitem[CL]{CL} O. Chodosh and C. Li, Generalized soap bubbles and the topology of manifolds with positive scalar curvature, Ann. of Math. 199 (2024), 707-740.
\bibitem[HKKZ]{HKKZ} S. Hirsch, D. Kazaras, M. Khuri, Y. Zhang, Spectral torical band inequalities and generalizations of the Schoen-Yau black hole existence theorem, International Mathematics Research Notices, Issue 4, (2024), 3139-3175.
\bibitem[LT]{LT} P. Li and L.F. Tam, Harmonic functions and the structure of complete manifolds, J. Differential Geom. 35 (1992) 359-383.
\bibitem[Li]{Li} P. Li, Geometric Analysis, Cambridge Stud. Adv. Math., 134, Cambridge University Press, 2012.
\bibitem[SY]{SY} R. Schoen and S. T. Yau, Existence of incompressible minimal surfaces and the topology of three-dimensional manifolds with nonnegative scalar curvature, Ann. Math. (2) 110 (1979), no. 1, 127-142.
\bibitem[Z]{Z} Jintian Zhu, Width estimate and doubly warped product, Trans. Amer. Math. Soc. 374 (2021), no. 2, 1497-1511.
\end{thebibliography}
\end{document}
